% !TEX program = xelatex
\documentclass[10pt,a4paper]{article}
\newcommand{\coursename}{Математический анализ — подготовка к экзамену}
\newcommand{\rightheader}{Теоретические вопросы}
\newcommand{\doctitle}{Теоретические вопросы}
% ==== Преамбула экзаменационных материалов (Theory / Proofs / Tasks) ====
% Поля уже, чем в конспектах; после каждого пункта оставлен небольшой
% вертикальный зазор для рукописных пометок.
% Перед этим файлом определяются: \documentclass и макросы
%   \coursename, \rightheader, \doctitle.
\usepackage{fontspec}
\setmainfont{cmunrm.otf}[
  BoldFont       = cmunbx.otf,
  ItalicFont     = cmunti.otf,
  BoldItalicFont = cmunbi.otf]
\usepackage{polyglossia}
\setdefaultlanguage{russian}
\setotherlanguage{english}

\usepackage{amsmath,amssymb,amsthm,mathtools}
\usepackage[a4paper,top=1.8cm,bottom=1.9cm,left=1.7cm,right=1.7cm]{geometry}
\usepackage{enumitem}
% ----- Менее сжатый текст и видимые отступы между абзацами -----
\linespread{1.28}
\setlength{\parskip}{0.55\baselineskip plus 2pt}
\setlist{itemsep=2pt,topsep=3pt,parsep=0pt}
\usepackage{graphicx}
\usepackage{array}
\usepackage{xcolor}
\definecolor{accent}{HTML}{1F6FEB}
\definecolor{teormin}{HTML}{C1121F}
\definecolor{rulecol}{HTML}{D0D7DE}
\usepackage[colorlinks=true,linkcolor=accent,urlcolor=accent,citecolor=accent]{hyperref}
\usepackage{titlesec}
\usepackage{fancyhdr}

% ----- Колонтитулы -----
\pagestyle{fancy}
\fancyhf{}
\fancyhead[L]{\small\itshape \coursename}
\fancyhead[R]{\small\itshape \rightheader}
\fancyfoot[C]{\thepage}
\renewcommand{\headrulewidth}{0.4pt}
\renewcommand{\headrule}{\hbox to\headwidth{\color{rulecol}\leaders\hrule height \headrulewidth\hfill}}

% ----- Заголовки разделов -----
\titleformat{\section}{\large\bfseries\color{accent}}{}{0pt}{}
\titlespacing*{\section}{0pt}{1.2ex}{0.6ex}
\titleformat{\subsection}{\bfseries}{}{0pt}{}
\titlespacing*{\subsection}{0pt}{1.0ex}{0.4ex}

% ----- Операторы и сокращения -----
\DeclareMathOperator{\tg}{tg}
\DeclareMathOperator{\ctg}{ctg}
\DeclareMathOperator{\arctg}{arctg}
\DeclareMathOperator{\arcctg}{arcctg}
\DeclareMathOperator{\sh}{sh}
\DeclareMathOperator{\ch}{ch}
\DeclareMathOperator{\Th}{th}
\DeclareMathOperator{\cth}{cth}
\DeclareMathOperator{\Const}{const}
\DeclareMathOperator{\sign}{sign}
\DeclareMathOperator{\Span}{span}
\newcommand{\R}{\mathbb{R}}
\newcommand{\C}{\mathbb{C}}
\newcommand{\N}{\mathbb{N}}
\newcommand{\Z}{\mathbb{Z}}
\newcommand{\Q}{\mathbb{Q}}
\newcommand{\dx}{\,dx}
\newcommand{\dt}{\,dt}
\newcommand{\du}{\,du}
\newcommand{\eps}{\varepsilon}
\renewcommand{\Re}{\operatorname{Re}}
\renewcommand{\Im}{\operatorname{Im}}
\renewcommand{\le}{\leqslant}
\renewcommand{\ge}{\geqslant}
\newcommand{\scal}[2]{\left( #1,\,#2\right)}
\DeclarePairedDelimiter{\abs}{\lvert}{\rvert}
\DeclarePairedDelimiter{\norm}{\lVert}{\rVert}

% ----- Окружения для доказательств (Proofs) -----
\theoremstyle{definition}
\newtheorem{definition}{Определение}
\newtheorem{example}{Пример}
\newtheorem{remark}{Замечание}
\theoremstyle{plain}
\newtheorem{theorem}{Теорема}
\newtheorem{lemma}{Лемма}
\renewcommand{\proofname}{Доказательство}

% ===== Хелперы для пунктов =====
% Регулируемый зазор для рукописных пометок (пустая строка после пункта):
\newcommand{\writespace}{\vspace{7.5mm}}

% Theory: \faq{номер}{вопрос}{ответ}
% номер теормина помечается красным значком !!! через \tm
\newcommand{\tm}{\textcolor{teormin}{\,!!!}}
\newcommand{\faq}[3]{%
  \par\smallskip\noindent
  \textbf{#1.}\enspace\textit{#2}\par
  \nobreak\vspace{0.3ex}\noindent #3\par
  \writespace}

% Proofs: \prf{номер}{заголовок} ... затем \stmt / \dok блоки
\newcommand{\prf}[2]{%
  \par\writespace\noindent
  {\large\bfseries\color{accent}#1.\enspace #2}\par\smallskip}
\newcommand{\stmt}{\par\noindent\textbf{Формулировка.}\enspace}
\newcommand{\dok}{\par\noindent\textbf{Доказательство.}\enspace}

% Tasks: \task{номер}{условие} ... затем \sol для решения
\newcommand{\task}[2]{%
  \par\writespace\noindent
  \textbf{\textcolor{accent}{Задача #1.}}\enspace #2\par\smallskip}
\newcommand{\sol}{\par\noindent\textbf{Решение.}\enspace}

\setlength{\parindent}{0pt}

\begin{document}
\begin{center}
  {\LARGE\bfseries \doctitle}\\[3pt]
  {\large\itshape \coursename}
\end{center}
\vspace{4pt}

\section{Часть I. Теоретический минимум \texorpdfstring{\textcolor{teormin}{(!!!)}}{(!!!)}}

\faq{1}{Табличные производные и интегралы.\tm}{%
$\int x^a\dx=\dfrac{x^{a+1}}{a+1}+C$ ($a\ne-1$),\;
$\int\dfrac{dx}{x}=\ln\abs{x}+C$,\;
$\int e^x\dx=e^x+C$,\;
$\int a^x\dx=\dfrac{a^x}{\ln a}+C$,\;
$\int\sin x\dx=-\cos x+C$,\;
$\int\cos x\dx=\sin x+C$,\;
$\int\dfrac{dx}{\cos^2x}=\tg x+C$,\;
$\int\dfrac{dx}{\sin^2x}=-\ctg x+C$,\;
$\int\dfrac{dx}{1+x^2}=\arctg x+C$,\;
$\int\dfrac{dx}{\sqrt{1-x^2}}=\arcsin x+C$,\;
$\int\dfrac{dx}{\sqrt{x^2\pm1}}=\ln\abs{x+\sqrt{x^2\pm1}}+C$.}

\faq{2}{Строгое определение первообразной для $f(x)$ на интервале $(a,b)$.\tm}{%
$F$ — первообразная для $f$ на $(a,b)$, если $F$ дифференцируема в каждой точке $(a,b)$ и
$F'(x)=f(x)$ для всех $x\in(a,b)$.}

\faq{3}{Теорема о замене переменной (условия и формула).\tm}{%
Если $x=\varphi(t)$ непрерывно дифференцируема, а $f$ имеет первообразную на множестве значений
$\varphi$, то $\displaystyle\int f(\varphi(t))\varphi'(t)\dt=\int f(x)\dx\big|_{x=\varphi(t)}$.}

\faq{4}{Производная от неопределённого интеграла и интеграл от дифференциала.\tm}{%
$\dfrac{d}{dx}\displaystyle\int f\dx=f(x)$,\quad $\displaystyle\int dF=F(x)+C$.}

\faq{5}{Определение неопределённого интеграла; геометрическая интерпретация.\tm}{%
$\int f\dx=\{F(x)+C\}$ — множество всех первообразных. Геометрически — семейство кривых, получаемых
вертикальным сдвигом: в каждой точке $x$ у них одинаковый наклон $f(x)$.}

\faq{6}{Теорема об общем виде первообразных на интервале.\tm}{%
Если $F$ — одна первообразная $f$ на интервале, то любая другая имеет вид $F(x)+C$; две
первообразные отличаются на постоянную.}

\faq{7}{Свойство линейности неопределённого интеграла.\tm}{%
$\int(\alpha f+\beta g)\dx=\alpha\int f\dx+\beta\int g\dx$.}

\faq{8}{Теорема о разложении правильной рациональной дроби на простейшие.\tm}{%
Множителю $(x-a)^k$ знаменателя отвечают слагаемые $\dfrac{A_1}{x-a}+\dots+\dfrac{A_k}{(x-a)^k}$;
множителю $(x^2+px+q)^m$ (неразложимому) — $\dfrac{M_1x+N_1}{x^2+px+q}+\dots+\dfrac{M_mx+N_m}{(x^2+px+q)^m}$.}

\faq{21}{Предел интегральных сумм при $\lambda(T)\to0$ на языке $\eps$–$\delta$.\tm}{%
Число $I$: $\forall\eps>0\ \exists\delta>0$ такое, что для любого разбиения с $\lambda(T)<\delta$ и
любой выборки $\xi$ выполнено $\abs{\sigma(f,T,\xi)-I}<\eps$.}

\faq{22}{Необходимое условие интегрируемости; контрпример к достаточности.\tm}{%
Интегрируемая функция \emph{ограничена}. Обратное неверно: функция Дирихле ограничена, но не
интегрируема по Риману.}

\faq{23}{Определения нижней $S_*(T)$ и верхней $S^*(T)$ сумм Дарбу.\tm}{%
$S_*(T)=\sum m_i\Delta x_i$, $S^*(T)=\sum M_i\Delta x_i$, где $m_i=\inf$, $M_i=\sup$ функции на
$[x_{i-1},x_i]$.}

\faq{24}{Нижний $\underline I$ и верхний $\overline I$ интегралы Дарбу; связь.\tm}{%
$\underline I=\sup_T S_*(T)$, $\overline I=\inf_T S^*(T)$; всегда $\underline I\le\overline I$.}

\faq{26}{Основной критерий интегрируемости через интегралы Дарбу.\tm}{%
Ограниченная $f$ интегрируема $\iff\underline I=\overline I$.}

\faq{29}{Определение равномерной непрерывности на множестве $X$.\tm}{%
$\forall\eps>0\ \exists\delta>0\ \forall x',x''\in X:\ \abs{x'-x''}<\delta\Rightarrow
\abs{f(x')-f(x'')}<\eps$ (один $\delta$ на всё $X$).}

\faq{30}{Теорема Кантора.\tm}{%
Функция, непрерывная на компакте (отрезке $[a,b]$), равномерно непрерывна на нём.}

\faq{31}{Критерий интегрируемости через $S^*(T)-S_*(T)<\eps$.\tm}{%
$f$ интегрируема $\iff\forall\eps>0\ \exists T:\ S^*(T)-S_*(T)<\eps$.}

\faq{35}{Точки разрыва первого и второго рода.\tm}{%
I рода: оба односторонних предела конечны (скачок или устранимый разрыв). II рода: хотя бы один
односторонний предел не существует или бесконечен.}

\faq{36}{Теорема об интегрируемости функций с конечным числом разрывов.\tm}{%
Ограниченная на $[a,b]$ функция с конечным числом точек разрыва интегрируема по Риману.}

\faq{38}{Линейность определённого интеграла.\tm}{%
$\int_a^b(\alpha f+\beta g)\dx=\alpha\int_a^b f\dx+\beta\int_a^b g\dx$.}

\faq{40}{Аддитивность по промежутку ($a<c<b$).\tm}{%
$\int_a^b f=\int_a^c f+\int_c^b f$.}

\faq{43}{Теорема об интегрировании неравенств.\tm}{%
Если $f(x)\le g(x)$ на $[a,b]$, то $\int_a^b f\le\int_a^b g$.}

\faq{44}{Неравенство для модуля интеграла.\tm}{%
$\abs[\big]{\int_a^b f}\le\int_a^b\abs{f}$ ($a<b$).}

\faq{46}{Первая теорема о среднем значении.\tm}{%
Для $f\in C[a,b]$ существует $\xi\in[a,b]$: $\int_a^b f=f(\xi)(b-a)$.}

\faq{48}{Теорема Вейерштрасса о точных гранях.\tm}{%
Непрерывная на отрезке функция достигает своих наибольшего и наименьшего значений.}

\faq{51}{Теорема Коши о промежуточном значении.\tm}{%
Непрерывная на отрезке функция принимает все значения между $m=\min$ и $M=\max$.}

\faq{52}{Определение $\Phi(x)=\int_a^x f(t)\dt$ и её область.\tm}{%
Для $f\in\mathcal R[a,b]$ функция $\Phi$ определена на всём $[a,b]$.}

\faq{53}{Теорема о непрерывности интеграла с переменным верхним пределом.\tm}{%
Если $f\in\mathcal R[a,b]$, то $\Phi(x)=\int_a^x f$ непрерывна (даже равномерно) на $[a,b]$.}

\faq{55}{Теорема Барроу.\tm}{%
Если $f$ непрерывна в точке $x$, то $\Phi$ дифференцируема в ней и $\Phi'(x)=f(x)$. Непрерывность
во всех точках не требуется — достаточно в исследуемой.}

\faq{61}{Теорема Ньютона–Лейбница.\tm}{%
Для $f\in C[a,b]$ и любой первообразной $F$: $\int_a^b f=F(b)-F(a)$.}

\faq{64}{Формула площади в полярных координатах.\tm}{%
$S=\dfrac12\int_\alpha^\beta\rho^2(\varphi)\,d\varphi$.}

\faq{67}{Определение длины дуги через вписанные ломаные.\tm}{%
Длина — точная верхняя грань (предел) длин вписанных ломаных при измельчении разбиения.}

\faq{69}{Три формулы длины дуги.\tm}{%
Декарт: $\int_a^b\sqrt{1+y'^2}\dx$; парам.: $\int\sqrt{x'^2+y'^2}\dt$; полярн.:
$\int\sqrt{\rho^2+\rho'^2}\,d\varphi$.}

\faq{70}{Объём тела вращения вокруг $Ox$.\tm}{%
$V=\pi\int_a^b f^2(x)\dx$.}

\faq{73}{Абсолютная и условная сходимость несобственного интеграла I рода.\tm}{%
Абсолютно — сходится $\int_a^{+\infty}\abs{f}$; условно — $\int f$ сходится, а $\int\abs{f}$
расходится.}

\faq{76}{Частичная сумма и сумма ряда; «ряд сходится».\tm}{%
$S_N=\sum_{n=1}^N a_n$. Ряд сходится, если существует конечный предел $\lim S_N=S$ (его сумма).}

\faq{77}{Необходимый признак сходимости; обратное.\tm}{%
Если ряд сходится, то $a_n\to0$. Обратное неверно: гармонический $\sum\frac1n$ расходится, хотя
$\frac1n\to0$.}

\faq{78}{Сумма эталонного геометрического ряда.\tm}{%
$\sum_{n=0}^\infty q^n=\dfrac1{1-q}$ при $\abs q<1$.}

\faq{79}{Критерий Коши для числового ряда.\tm}{%
$\forall\eps>0\ \exists N\ \forall n>N\ \forall p:\ \abs{a_{n+1}+\dots+a_{n+p}}<\eps$.}

\faq{80}{Остаток $R_K$ ряда; связь со сходимостью.\tm}{%
$R_K=\sum_{n=K+1}^\infty a_n$. Ряд сходится $\iff R_K\to0$.}

\faq{86}{Признак сравнения в форме неравенств.\tm}{%
Пусть $0\le a_n\le b_n$. Из сходимости $\sum b_n$ следует сходимость $\sum a_n$; из расходимости
$\sum a_n$ — расходимость $\sum b_n$.}

\faq{87}{Признак сравнения в предельной форме.\tm}{%
Если $a_n,b_n>0$ и $\frac{a_n}{b_n}\to L\in(0,\infty)$, ряды сходятся или расходятся одновременно.}

\faq{88}{Признак Д'Аламбера (предельная форма).\tm}{%
$\lim\frac{a_{n+1}}{a_n}=q$: при $q<1$ сходится, $q>1$ расходится, $q=1$ — не даёт ответа.}

\faq{89}{Обобщённый гармонический ряд.\tm}{%
$\sum\frac1{n^p}$ сходится при $p>1$ и расходится при $p\le1$.}

\faq{90}{Радикальный признак Коши (предельная форма).\tm}{%
$\lim\sqrt[n]{a_n}=q$: при $q<1$ сходится, $q>1$ расходится, $q=1$ — не даёт ответа.}

\faq{92}{Определение абсолютной сходимости.\tm}{%
$\sum a_n$ сходится абсолютно, если сходится $\sum\abs{a_n}$.}

\faq{95}{Определение условной сходимости.\tm}{%
Ряд сходится, но $\sum\abs{a_n}$ расходится.}

\faq{98}{Признак Лейбница.\tm}{%
Если $c_n>0$ монотонно убывают и $c_n\to0$, то $\sum(-1)^{n-1}c_n$ сходится.}

\faq{110}{Поточечная сходимость $f_n\to f$ (в кванторах).\tm}{%
$\forall x\in X\ \forall\eps>0\ \exists N(\eps,x)\ \forall n>N:\ \abs{f_n(x)-f(x)}<\eps$.}

\faq{111}{Равномерная сходимость $f_n\rightrightarrows f$; отличие.\tm}{%
$\forall\eps>0\ \exists N(\eps)\ \forall n>N\ \forall x\in X:\ \abs{f_n(x)-f(x)}<\eps$. Отличие:
$N$ не зависит от $x$ ($\exists N$ стоит \emph{до} $\forall x$).}

\faq{115}{Критерий Коши равномерной сходимости последовательности.\tm}{%
$\forall\eps>0\ \exists N\ \forall n>N\ \forall p\ \forall x:\ \abs{f_{n+p}(x)-f_n(x)}<\eps$.}

\faq{116}{Критерий Коши равномерной сходимости ряда.\tm}{%
$\forall\eps>0\ \exists N\ \forall n>N\ \forall p\ \forall x:\
\abs[\big]{\sum_{k=n+1}^{n+p}u_k(x)}<\eps$.}

\faq{119}{Мажорантный признак Вейерштрасса.\tm}{%
Если $\abs{u_n(x)}\le a_n$ для всех $x$ и $\sum a_n$ сходится, то $\sum u_n$ сходится равномерно (и
абсолютно).}

\faq{121}{Теорема о непрерывности суммы ряда.\tm}{%
Если $u_n\in C(X)$ и $\sum u_n\rightrightarrows S$, то $S\in C(X)$.}

\faq{124}{Теорема о почленном интегрировании ряда.\tm}{%
Если $u_n\in C[a,b]$ и $\sum u_n\rightrightarrows S$, то $\int_a^b\sum u_n=\sum\int_a^b u_n$;
промежуток — конечный отрезок.}

\faq{125}{Теорема о почленном дифференцировании ряда.\tm}{%
Если $u_n\in C^1$, $\sum u_n$ сходится в точке и $\sum u_n'\rightrightarrows$ равномерно, то
$(\sum u_n)'=\sum u_n'$. Равномерной должна быть сходимость \emph{ряда производных}.}

\faq{127}{Определение степенного ряда по степеням $(x-x_0)$.\tm}{%
$\sum_{n=0}^\infty a_n(x-x_0)^n$.}

\faq{128}{Первая теорема Абеля.\tm}{%
Если ряд сходится в $x_1\ne x_0$, то он абсолютно сходится при всех $x$ с $\abs{x-x_0}<\abs{x_1-x_0}$.}

\faq{129}{Геометрическая форма области сходимости.\tm}{%
Интервал $(x_0-R,x_0+R)$, симметричный относительно центра (возможно, с одним или обоими концами).}

\faq{130}{Формула Даламбера для радиуса сходимости.\tm}{%
$R=\lim\abs[\big]{\dfrac{a_n}{a_{n+1}}}$, если этот предел существует.}

\faq{131}{Теорема Коши–Адамара.\tm}{%
$\dfrac1R=\varlimsup\limits_{n\to\infty}\sqrt[n]{\abs{a_n}}$.}

\faq{132}{Диагноз при $R=0$ и $R=+\infty$.\tm}{%
$R=0$: ряд сходится только в центре $x_0$; $R=+\infty$: сходится на всей $\R$.}

\faq{133}{Сходимость на вложенном отрезке $[\alpha,\beta]\subset(-R,R)$.\tm}{%
На любом таком отрезке ряд сходится равномерно.}

\faq{134}{Теоремы о свойствах степенного ряда внутри интервала.\tm}{%
Сумма непрерывна; ряд можно почленно интегрировать и (бесконечно) дифференцировать на
$(x_0-R,x_0+R)$.}

\faq{135}{Меняется ли $R$ при дифференцировании/интегрировании?\tm}{%
Нет, радиус сходимости сохраняется.}

\faq{139}{Формула коэффициентов Тейлора.\tm}{%
$a_n=\dfrac{f^{(n)}(x_0)}{n!}$.}

\faq{140}{Формальный вид ряда Тейлора.\tm}{%
$\sum_{n=0}^\infty\dfrac{f^{(n)}(x_0)}{n!}(x-x_0)^n$.}

\faq{145}{Четыре аксиомы скалярного произведения.\tm}{%
1) $(f,g)=(g,f)$; 2) $(\alpha f,g)=\alpha(f,g)$; 3) $(f_1+f_2,g)=(f_1,g)+(f_2,g)$;
4) $(f,f)\ge0$, причём $=0\iff f\equiv0$.}

\faq{146}{Неравенство Коши–Буняковского в $L_2[a,b]$.\tm}{%
$\abs[\big]{\int_a^b fg}\le\sqrt{\int_a^b f^2}\cdot\sqrt{\int_a^b g^2}$.}

\faq{148}{Ортогональная и ортонормированная системы.\tm}{%
Ортогональна: $(\varphi_n,\varphi_m)=0$ при $n\ne m$. Ортонормирована: ещё и $\norm{\varphi_n}=1$.}

\faq{149}{Формула коэффициентов Фурье по ортогональной системе.\tm}{%
$c_n=\dfrac{(f,\varphi_n)}{\norm{\varphi_n}^2}$.}

\faq{155}{Признак Дини.\tm}{%
Если $\displaystyle\int_0^\pi\frac{\abs{f(x_0+t)+f(x_0-t)-2S}}{t}\dt<\infty$, то ряд Фурье в точке
$x_0$ сходится к $S$.}

\faq{157}{Сумма ряда Фурье в точке разрыва I рода.\tm}{%
Ряд сходится к полусумме односторонних пределов $\dfrac{f(x_0-0)+f(x_0+0)}{2}$.}

\faq{159}{Аксиома эрмитовости комплексного скалярного произведения.\tm}{%
$(f,g)=\overline{(g,f)}$ (косая симметрия); тогда $(f,f)=\int\abs{f}^2\ge0$ вещественно.}

\faq{160}{Комплексный коэффициент Фурье $c_n$.\tm}{%
$c_n=\dfrac1{2\pi}\int_{-\pi}^\pi f(x)e^{-inx}\dx$.}

\section{Часть II. Дополнительные вопросы}

\faq{9}{Канонический вид элементарных дробей I–IV типов.}{%
I: $\dfrac{A}{x-a}$,\; II: $\dfrac{A}{(x-a)^k}$,\; III: $\dfrac{Mx+N}{x^2+px+q}$,\;
IV: $\dfrac{Mx+N}{(x^2+px+q)^k}$ ($p^2-4q<0$, $k\ge2$).}

\faq{10}{Интеграл дроби III типа; полный квадрат в $x^2+px+q$.}{%
$x^2+px+q=\bigl(x+\tfrac p2\bigr)^2+\bigl(q-\tfrac{p^2}4\bigr)$; интеграл распадается на
логарифмическую ($\ln$) и арктангенсную ($\arctg$) части.}

\faq{11}{Универсальная тригонометрическая подстановка.}{%
$t=\tg\tfrac x2$: $\sin x=\dfrac{2t}{1+t^2}$, $\cos x=\dfrac{1-t^2}{1+t^2}$,
$dx=\dfrac{2\,dt}{1+t^2}$.}

\faq{12}{Уравнение-шаблон метода Остроградского.}{%
$\int\frac PQ\dx=\frac{P_1}{Q_1}+\int\frac{P_2}{Q_2}\dx$, где $Q_1=\gcd(Q,Q')$, $Q_2=Q/Q_1$,
$\deg P_1<\deg Q_1$, $\deg P_2<\deg Q_2$.}

\faq{13}{Когда метод Остроградского целесообразен.}{%
Когда знаменатель имеет корни высокой кратности — метод сразу выделяет рациональную часть без
громоздкого разложения на простейшие.}

\faq{14}{Дробно-линейная подстановка.}{%
Для $R\bigl(x,\sqrt[n]{\tfrac{ax+b}{cx+d}}\bigr)$: подстановка $t^n=\dfrac{ax+b}{cx+d}$.}

\faq{15}{Замена для радикала $\sqrt{x^2+q}$.}{%
$x=\sqrt q\,\sh t$, $dx=\sqrt q\,\ch t\dt$, $\sqrt{x^2+q}=\sqrt q\,\ch t$ (гиперболическая
подстановка).}

\faq{16}{Три подстановки Эйлера для $\sqrt{ax^2+bx+c}$.}{%
1) $a>0$: $\sqrt{ax^2+bx+c}=\pm\sqrt a\,x+t$; 2) $c>0$: $\sqrt{\dots}=xt\pm\sqrt c$;
3) есть вещественные корни $x_1,x_2$: $\sqrt{\dots}=t(x-x_1)$.}

\faq{17}{Дифференциальный бином Чебышёва.}{%
$\int x^m(a+bx^n)^p\dx$, где $a,b\in\R$, а $m,n,p\in\Q$.}

\faq{18}{Три условия интегрируемости (Чебышёв).}{%
Берётся в элементарных функциях $\iff$ (1) $p\in\Z$; (2) $\tfrac{m+1}{n}\in\Z$; (3)
$\tfrac{m+1}{n}+p\in\Z$.}

\faq{19}{Подстановки по случаям Чебышёва.}{%
(1) $x=t^N$, $N$ — НОК знаменателей $m,n$; (2) $a+bx^n=t^s$, $s$ — знаменатель $p$; (3)
$ax^{-n}+b=t^s$.}

\faq{20}{Определение интегральной суммы Римана.}{%
$\sigma(f,T,\{\xi_i\})=\sum_{i=1}^n f(\xi_i)\Delta x_i$, где $T:\,a=x_0<\dots<x_n=b$,
$\xi_i\in[x_{i-1},x_i]$.}

\faq{25}{Колебание $\omega_i$ функции.}{%
$\omega_i=M_i-m_i=\sup-\inf$ на $[x_{i-1},x_i]=\sup_{x',x''}\abs{f(x')-f(x'')}$.}

\faq{27}{Критерий интегрируемости через колебания.}{%
$f$ интегрируема $\iff\lim_{\lambda(T)\to0}\sum\omega_i\Delta x_i=0$.}

\faq{28}{Геометрический смысл $S^*(T)-S_*(T)$.}{%
Суммарная площадь «зазора» между описанной и вписанной ступенчатыми фигурами.}

\faq{32}{Монотонная функция; обязана ли быть непрерывной?}{%
Неубывающая: $x_1<x_2\Rightarrow f(x_1)\le f(x_2)$ (невозрастающая — наоборот). Непрерывной быть
не обязана.}

\faq{33}{Ограниченность из монотонности.}{%
Неубывающая на $[a,b]$ ограничена: $f(a)\le f(x)\le f(b)$; $\inf=f(a)$, $\sup=f(b)$.}

\faq{34}{Пример монотонной функции со счётным числом разрывов.}{%
Функция скачков $f(x)=\sum_{r_n<x}2^{-n}$ по всем рациональным $r_n\in[0,1]$: возрастает, разрывна
в каждой рациональной точке.}

\faq{37}{$f(x)=1/x$ на $(0,1]$ — интегрируема ли?}{%
Нет: функция неограничена, а ограниченность — необходимое условие интегрируемости по Риману.}

\faq{39}{Свойство сохранения знака.}{%
Если $f\ge0$ на $[a,b]$, то $\int_a^b f\ge0$.}

\faq{41}{Интеграл при $a=b$ и $a>b$.}{%
$\int_a^a f=0$, $\int_a^b f=-\int_b^a f$; формальная аддитивность сохраняется.}

\faq{42}{Физический смысл аддитивности.}{%
Масса стержня $[a,b]$ равна сумме масс $[a,c]$ и $[c,b]$ (аналогично — работа, заряд).}

\faq{45}{Существование $\int\abs{f}$; обратное.}{%
Если $f\in\mathcal R[a,b]$, то и $\abs f\in\mathcal R[a,b]$. Обратное неверно: $\abs f$ может быть
интегрируема, а $f$ — нет (напр.\ $f=\pm1$ по рациональности $x$).}

\faq{47}{Среднее значение функции.}{%
$\mu=\dfrac1{b-a}\int_a^b f(x)\dx$.}

\faq{49}{Обобщённая теорема о среднем.}{%
Для $f\in C$, $g\in\mathcal R[a,b]$ знакопостоянной: $\int_a^b fg=f(\xi)\int_a^b g$, $\xi\in[a,b]$.}

\faq{50}{Ограничение на $g$ в обобщённой теореме.}{%
$g$ должна сохранять знак на $[a,b]$; без этого теорема перестаёт быть верной.}

\faq{54}{Непрерывность $\Phi$ в точке разрыва I рода $f$.}{%
$\Phi$ остаётся непрерывной: интеграл с переменным пределом непрерывен при любой интегрируемой
$f$ (скачок «сглаживается»).}

\faq{56}{Производная интеграла с двумя переменными пределами.}{%
$\dfrac{d}{dx}\int_{\psi(x)}^{\varphi(x)}f(t)\dt=f(\varphi(x))\varphi'(x)-f(\psi(x))\psi'(x)$.}

\faq{57}{$\Phi(x)=\int_0^x\abs t\dt$ в точке $x=0$.}{%
$\Phi'(0)=\abs0=0$ — подынтегральная функция непрерывна в нуле, по Барроу $\Phi'(0)=f(0)=0$.}

\faq{58}{Условия применимости Лопиталя к интегралам.}{%
Числитель и знаменатель $\to0$ (или $\to\infty$), дифференцируемы, производная знаменателя $\ne0$
вблизи $x_0$, и существует предел отношения производных.}

\faq{59}{Теорема о существовании первообразной.}{%
У всякой непрерывной на $[a,b]$ функции есть первообразная $\Phi(x)=\int_a^x f$ — прямое следствие
теоремы Барроу.}

\faq{60}{$\displaystyle\lim_{x\to0+}\frac{\int_0^x\sin(t^2)\dt}{x^3}$.}{%
$\sin(t^2)\sim t^2$, $\int_0^x t^2\dt=\tfrac{x^3}3$, предел равен $\tfrac13$.}

\faq{62}{Разница определённого и неопределённого интегралов.}{%
Определённый интеграл — \emph{число}; неопределённый — \emph{семейство функций} (множество
первообразных).}

\faq{63}{Н–Л для $f(x)=\sign x$ на $[-1,1]$.}{%
Неприменима: $\sign x$ не имеет первообразной на $[-1,1]$ (разрыв в нуле), хотя сам интеграл
существует и равен $0$.}

\faq{65}{Условия формулы площади в полярных координатах.}{%
$\rho(\varphi)\ge0$ непрерывна на $[\alpha,\beta]$, и $\beta-\alpha\le2\pi$ (сектор не
перекрывается).}

\faq{66}{Площадь кругового сектора.}{%
$S=\dfrac12 R^2\Delta\varphi$ — базовый элемент аппроксимации.}

\faq{68}{Спрямляемая кривая; достаточное условие.}{%
Кривая конечной длины ($\sup$ длин ломаных $<\infty$). Достаточно гладкости (непрерывная
производная).}

\faq{71}{Объём вращения вокруг $Oy$.}{%
$V=2\pi\int_a^b x\,f(x)\dx$ (метод оболочек, $a\ge0$).}

\faq{72}{Элемент объёма $dV$: диски и оболочки.}{%
Диск: $dV=\pi f^2(x)\dx$ (тонкий цилиндр радиуса $f$). Оболочка: $dV=2\pi x\,f(x)\dx$ (тонкая труба
радиуса $x$).}

\faq{74}{Признак Абеля для несобственного интеграла.}{%
$\int_a^{+\infty}fg$ сходится, если $\int f$ сходится, а $g$ монотонна и ограничена.}

\faq{75}{Интеграл Дирихле.}{%
$\int_0^{+\infty}\dfrac{\sin x}{x}\dx=\dfrac\pi2$; сходится \emph{условно} (не абсолютно).}

\faq{81}{Отрицание критерия Коши (расходимость).}{%
$\exists\eps>0\ \forall N\ \exists n>N\ \exists p:\ \abs{a_{n+1}+\dots+a_{n+p}}\ge\eps$.}

\faq{82}{Линейность сходящихся рядов.}{%
Если $\sum a_n=A$, $\sum b_n=B$, то $\sum(\alpha a_n+\beta b_n)=\alpha A+\beta B$.}

\faq{83}{$\sum a_n$ сходится, $\sum b_n$ расходится — что с $\sum(a_n+b_n)$?}{%
Расходится.}

\faq{84}{Отбрасывание первых $10^6$ членов.}{%
Факт сходимости не меняется; сумма меняется (на сумму отброшенных членов).}

\faq{85}{Критерий сходимости ряда с неотрицательными членами.}{%
Ряд с $a_n\ge0$ сходится $\iff$ его частичные суммы ограничены сверху.}

\faq{91}{Что «сильнее» — Д'Аламбер или Коши?}{%
Коши сильнее: если $\exists\lim\frac{a_{n+1}}{a_n}=q$, то $\exists\lim\sqrt[n]{a_n}=q$, но не
наоборот.}

\faq{93}{Перестановка членов абсолютно сходящегося ряда (Дирихле).}{%
Любая перестановка абсолютно сходящегося ряда сходится к той же сумме.}

\faq{94}{$\sum a_n$ абс.\ $+\sum b_n$ усл.\ — что с $\sum(a_n+b_n)$?}{%
Сходится условно.}

\faq{96}{Знакочередующийся гармонический ряд.}{%
$\sum_{n=1}^\infty\dfrac{(-1)^{n-1}}{n}=\ln2$.}

\faq{97}{Необходимое условие для знакопеременного ряда.}{%
$a_n\to0$: $\forall\eps>0\ \exists N\ \forall n>N:\ \abs{a_n}<\eps$.}

\faq{99}{Оценка остатка ряда Лейбница.}{%
$\abs{R_N}\le c_{N+1}$ — не превосходит модуля первого отброшенного члена.}

\faq{100}{Знак остатка $R_N$ ряда Лейбница.}{%
Совпадает со знаком первого отброшенного члена.}

\faq{101}{Формула преобразования Абеля.}{%
$\sum_{k=1}^n a_k b_k=a_n B_n-\sum_{k=1}^{n-1}(a_{k+1}-a_k)B_k$, где $B_k=\sum_{j=1}^k b_j$.}

\faq{102}{Лемма Абеля об оценке частичных сумм.}{%
Если $\{a_k\}$ монотонна и $\abs{B_k}\le M$, то $\abs[\big]{\sum a_k b_k}\le M(\abs{a_1}+2\abs{a_n})$.}

\faq{103}{Аналогия преобразования Абеля с интегрированием по частям.}{%
Разность $\Delta a_k$ ↔ дифференциал $da$, частичная сумма $B_k$ ↔ интеграл $\int b$; формула
$\sum a\,\Delta B=aB-\sum B\,\Delta a$ повторяет $\int u\,dv=uv-\int v\,du$.}

\faq{104}{Признак Дирихле для числового ряда.}{%
$\sum a_n b_n$ сходится, если частичные суммы $B_n=\sum b_k$ ограничены, а $a_n\to0$ монотонно.}

\faq{105}{Сходимость эталонного $\sum\frac{\sin n}{n^p}$.}{%
Сходится при $p>0$ (Дирихле); абсолютно при $p>1$; при $0<p\le1$ — лишь условно.}

\faq{106}{Замкнутая сумма $\sum_{k=1}^n\sin kx$.}{%
$\sum_{k=1}^n\sin kx=\dfrac{\sin\frac{nx}2\,\sin\frac{(n+1)x}2}{\sin\frac x2}$; при $x\ne2\pi m$
ограничена по $n$.}

\faq{107}{Признак Абеля для числового ряда.}{%
$\sum a_n b_n$ сходится, если $\sum b_n$ сходится, а $a_n$ монотонна и ограничена.}

\faq{108}{Различие признаков Дирихле и Абеля.}{%
Дирихле: $a_n\to0$ монотонно $+$ ограниченные суммы $b$. Абель: $a_n$ монотонна и ограничена (к
нулю не обязана) $+$ сходящийся $\sum b_n$.}

\faq{109}{$\sum b_n$ сх., $a_n$ монот.\ огранич.\ — сходится ли $\sum\abs{a_n b_n}$?}{%
Не обязан: признак Абеля даёт лишь (возможно условную) сходимость $\sum a_n b_n$, не абсолютную.}

\faq{112}{Геометрическая интерпретация равномерной сходимости.}{%
Начиная с номера $N$, график $f_n$ целиком лежит в $\eps$-полосе $\{\abs{y-f(x)}<\eps\}$ вокруг
$f$.}

\faq{113}{Супремум-критерий равномерной сходимости.}{%
$f_n\rightrightarrows f\iff\rho_n=\sup_X\abs{f_n-f}\to0$.}

\faq{114}{Поточечно к непрерывной $\Rightarrow$ равномерно?}{%
Нет. Контрпример «бегущего горба»: $f_n\to0$ (предел непрерывен), но $\rho_n\not\to0$.}

\faq{117}{Отличие критерия Коши от супремум-критерия.}{%
Критерий Коши \emph{не требует} знания предельной функции $f$ (внутренняя характеристика);
супремум-критерий — требует.}

\faq{118}{Равномерная сходимость ряда через остатки.}{%
$\sum u_n\rightrightarrows S\iff\sup_X\abs{R_n(x)}\to0$, где $R_n=\sum_{k>n}u_k$.}

\faq{120}{Равномерно $\Rightarrow$ абсолютно?}{%
Нет. Контрпример: $u_n(x)=\frac{(-1)^n}{n}$ (константы) — ряд сходится равномерно, но $\sum\frac1n$
расходится.}

\faq{122}{Перестановка предела и суммы.}{%
Если $\sum u_n\rightrightarrows S$ и $\lim_{x\to x_0}u_n(x)=b_n$, то
$\lim_{x\to x_0}\sum u_n(x)=\sum b_n$.}

\faq{123}{Пример: предел непрерывных — разрывная функция.}{%
$f_n(x)=x^n$ на $[0,1]$: предел равен $0$ при $x<1$ и $1$ при $x=1$ (разрывен).}

\faq{126}{Уравнение почленного дифференцирования.}{%
$\dfrac{d}{dx}\sum_{n=1}^\infty u_n(x)=\sum_{n=1}^\infty u_n'(x)$.}

\faq{136}{Аналитичность в точке (по Тао).}{%
$f$ аналитична в $x_0$, если в некоторой окрестности $f(x)=\sum a_n(x-x_0)^n$ — сходящийся
степенной ряд.}

\faq{137}{Иерархия классов $C^0,C^1,C^\infty,C^\omega$.}{%
$C^\omega\subset C^\infty\subset\dots\subset C^1\subset C^0$; глубже всех вложен $C^\omega$
(аналитические функции).}

\faq{138}{Необходимое условие аналитичности.}{%
$f\in C^\infty$ (бесконечная дифференцируемость).}

\faq{141}{Когда ряд Тейлора называется рядом Маклорена.}{%
При центре разложения $x_0=0$.}

\faq{142}{Остаточный член $R_n$ в форме Лагранжа и Коши.}{%
Лагранж: $R_n=\dfrac{f^{(n+1)}(\xi)}{(n+1)!}(x-x_0)^{n+1}$; Коши:
$R_n=\dfrac{f^{(n+1)}(\xi)}{n!}(x-\xi)^n(x-x_0)$.}

\faq{143}{Критерий сходимости ряда Тейлора к функции.}{%
$R_n(x)\to0$ при $n\to\infty$.}

\faq{144}{Достаточное условие аналитичности.}{%
Если в окрестности $\abs{f^{(n)}(x)}\le C\cdot M^n$ для всех $n$, то $f$ аналитична.}

\faq{147}{Проблема положительной определённости; фактор-пространство.}{%
$\int f^2=0$ не влечёт $f\equiv0$ ($f$ может быть $\ne0$ на множестве меры нуль). Исправляется
фактор-пространством — склейкой функций, совпадающих почти всюду.}

\faq{150}{Геометрический смысл Бесселя и Парсеваля.}{%
Бесконечномерная теорема Пифагора: $\norm{f}^2\ge\sum c_n^2\norm{\varphi_n}^2$ (Бессель), с
равенством (Парсеваль) при полноте системы.}

\faq{151}{Интегральное представление $S_N(x)$.}{%
$S_N(x)=\dfrac1\pi\int_{-\pi}^\pi f(x+t)D_N(t)\dt$ (свёртка с ядром Дирихле).}

\faq{152}{Формула ядра Дирихле и три свойства.}{%
$D_N(t)=\dfrac{\sin\bigl((N+\frac12)t\bigr)}{2\sin\frac t2}$; чётно, нормировано
($\frac1\pi\int_{-\pi}^\pi D_N=1$), осциллирует.}

\faq{153}{Лемма Римана–Лебега.}{%
Для интегрируемой $f$: $\int f(x)\sin(\lambda x)\dx\to0$ и $\int f(x)\cos(\lambda x)\dx\to0$ при
$\lambda\to\infty$.}

\faq{154}{Условие Липшица порядка $\alpha$.}{%
$\abs{f(x)-f(x_0)}\le L\abs{x-x_0}^\alpha$, $0<\alpha\le1$. Обычная дифференцируемость даёт
$\alpha=1$.}

\faq{156}{Сходится ли ряд Фурье непрерывной функции всюду?}{%
Нет, не обязательно: существуют непрерывные функции с расходящимся в точке рядом Фурье.}

\faq{158}{Ряд Лейбница для $\pi$.}{%
$\dfrac\pi4=1-\dfrac13+\dfrac15-\dfrac17+\dots$}

\faq{161}{Связь $a_n,b_n$ с $c_n,c_{-n}$.}{%
$c_n=\dfrac{a_n-ib_n}2$, $c_{-n}=\dfrac{a_n+ib_n}2$; обратно $a_n=c_n+c_{-n}$, $b_n=i(c_n-c_{-n})$.}

\end{document}
